\documentclass[11pt,a4paper]{article}
\usepackage[utf8]{inputenc}
\usepackage[T1]{fontenc}
\usepackage{amsmath,amssymb,amsthm}
\usepackage{mathtools}
\usepackage{mathrsfs}
\usepackage[colorlinks=true,linkcolor=blue,citecolor=blue,urlcolor=blue]{hyperref}
\usepackage{booktabs}

\newtheorem{theorem}{Theorem}[section]
\newtheorem{lemma}[theorem]{Lemma}
\newtheorem{proposition}[theorem]{Proposition}
\newtheorem{corollary}[theorem]{Corollary}
\newtheorem{definition}[theorem]{Definition}
\theoremstyle{remark}
\newtheorem{remark}[theorem]{Remark}

\title{\bfseries Research Programme III\\[4pt]
\large The Categorical Explicit Formula and the Limit to \(\zeta(2s)\)}
\author{Victor Geere}
\date{\today}

\begin{document}

\maketitle

\begin{abstract}
Two earlier research programmes attempted to link a greedy harmonic sieve to the Riemann Hypothesis; both failed because the sieve destroys the multiplicative structure essential to the zeta function.
The present programme abandons sieves entirely and works directly with the finite‑level categorical Dirichlet polynomials
\[
\Psi_k(s)=\sum_{j}\chi_{j,k}\,m_{j,k}^{-s},
\]
which arise from the tower \(\mathcal{C}_k=\mathrm{SU}(2)_k\) of modular tensor categories.
Each \(\Psi_k(s)\) satisfies a Riemann‑zeta‑type functional equation (the Quantum‑Egyptian Functional Equation) and, after normalisation, converges to \(\zeta(2s)\) as \(k\to\infty\).
The zeros of the completed function \(\Lambda_k(s)\) are finite in number and lie symmetrically about the critical line \(\Re(s)=1/2\).
We propose to study these zeros via an explicit formula that relates them to the underlying theta function and Galois action.
The programme is a purely spectral approach: the modular \(S\)-matrix, its Galois symmetries, and the twisted theta function are the central objects.
Its ultimate goal is to prove that, in the limit, all zeros lie on the critical line, thereby establishing the Riemann Hypothesis for \(\zeta(s)\).
\end{abstract}

\section{Introduction and Motivation}

\subsection{Lessons from Earlier Attempts}

\textbf{Programme I} (the MDS Unified System) split the von Mangoldt measure using a greedy harmonic selection on prime powers.
Its analysis revealed that the critical thresholds simplify to \(\Theta_g=\pi/g\), that the budget dynamics are deterministic non‑contractive shifts, and that the resulting Dirichlet series has no Euler product.
The companion paper \cite{greedy-sieve} proved rigorously that the greedy harmonic sieve is \emph{not} an arithmetic sieve: its indicator function is not multiplicative, its Dirichlet series lacks an Euler product, and its analytic properties are unrelated to those of \(\zeta(s)\).

\textbf{Programme II} attempted to port the sieve to modular tensor categories by sieving the Egyptian denominators \(m_i=D^2/d_i^2\).
The categorical MDS failed because fusion rings lack unique factorisation, and the integer MDS on the \(m_i\) again produced a sub‑series with no intrinsic functional equation.
The sieve systematically destroyed the Galois‑sign multiplicativity that makes the categorical Dirichlet series interesting.

The positive outcome of both investigations is the robustness of the \textbf{Quantum‑Egyptian Functional Equation} (QEFE) \cite{qefe}.
For every integral modular tensor category, the series \(\Psi_{\mathcal{C},\chi}(s)=\sum_i\chi(m_i)m_i^{-s}\) satisfies a Riemann‑zeta‑type functional equation, proved via the modular transformation of a twisted theta function.
For the tower \(\mathrm{SU}(2)_k\), numerical evidence indicates that the normalised series \(\widehat{\Psi}_k(s)=\Psi_k(s)/\Psi_k(1/2)\) converges to \(\zeta(2s)\).
This convergence is \emph{not} mediated by a sieve; it is a direct limit of finite Dirichlet polynomials with exact symmetries.

\subsection{The Present Programme}

We propose to study the finite‑level Dirichlet polynomials \(\Psi_k(s)\) as approximations to \(\zeta(2s)\), using the spectral data provided by the QEFE.
The zeros of each completed function \(\Lambda_k(s)\) are finite, symmetric about \(\Re(s)=1/2\), and can be analysed through an explicit formula linking them to the coefficients and the theta function.
The ultimate goal is to understand whether the limit of these zero distributions forces all zeros onto the critical line, thereby proving the Riemann Hypothesis for \(\zeta(2s)\) and hence for \(\zeta(s)\).

\section{Preliminaries}

All definitions and results in this section are taken from the source documents \cite{greedy-continuation,greedy-sieve,qefe}, which should be consulted for full proofs.

\subsection{The Tower \(\mathrm{SU}(2)_k\) and Egyptian Denominators}

For a fixed level \(k\in\mathbb{Z}_{\ge 1}\), the modular tensor category \(\mathcal{C}_k=\mathrm{SU}(2)_k\) has simple objects labelled by spins \(j=0,\frac12,1,\dots,\frac{k}{2}\).
Their quantum dimensions are
\begin{equation}\label{eq:qdim}
d_{j,k}= \frac{\sin\!\bigl((2j+1)\frac{\pi}{k+2}\bigr)}{\sin\!\bigl(\frac{\pi}{k+2}\bigr)},
\end{equation}
and the total quantum dimension is
\begin{equation}\label{eq:D2}
D_k^2 = \sum_{j} d_{j,k}^2 = \frac{k+2}{4\sin^2\!\bigl(\frac{\pi}{k+2}\bigr)}.
\end{equation}
The \textbf{Egyptian denominators} are the integers
\begin{equation}\label{eq:m}
m_{j,k} = \frac{D_k^2}{d_{j,k}^2}
        = \frac{k+2}{4\sin^2\!\bigl((2j+1)\frac{\pi}{k+2}\bigr)}.
\end{equation}
For large \(k\) one has the asymptotic expansion
\begin{equation}\label{eq:masympt}
m_{j,k} \sim \frac{(k+2)^3}{3\pi^2(2j+1)^2} \qquad (k\to\infty,\;j\text{ fixed}).
\end{equation}
All \(m_{j,k}\) are integers.

\subsection{Galois Action and Character \(\chi\)}

The modular \(S\)-matrix of \(\mathcal{C}_k\) has entries
\[
S_{j,j'} = \sqrt{\frac{2}{k+2}}\, \sin\!\Bigl((2j+1)(2j'+1)\frac{\pi}{k+2}\Bigr).
\]
Normalising by \(D_k\) gives the unitary matrix \(s_{j,j'}=S_{j,j'}/D_k\).
The Galois group \(\mathrm{Gal}(\mathbb{Q}(\zeta_{4(k+2)})/\mathbb{Q})\) acts on the entries: for an automorphism \(\sigma\) sending \(\zeta_{k+2}\mapsto\zeta_{k+2}^\ell\) with \((\ell,k+2)=1\), the matrix \(s\) is mapped to a signed permutation matrix (Huang~\cite{huang2019}).
The signs
\[
\chi_{j,k}=\chi(m_{j,k})\in\{\pm1\}
\]
are defined as the sign attached to the object \(j\) under this action.
They form a multiplicative character in the sense that they respect the tensor product decomposition (they arise from a homomorphism of the Grothendieck ring to \(\{\pm1\}\)).

\subsection{The Twisted Theta Function}

From the Egyptian denominators and the Galois signs one builds the theta function
\begin{equation}\label{eq:theta}
\Theta_k(\tau) = \sum_{j} \chi_{j,k}\, e^{-\pi (m_{j,k}/D_k)^2 \tau}, \qquad \tau\in\mathbb{H}.
\end{equation}
Using the modular transformation of the \(S\)-matrix one proves a modular relation
\begin{equation}\label{eq:modrel}
\Theta_k(-1/\tau) = \varepsilon_k\, \tau^{a_k}\, \overline{\Theta_k(\tau)},
\end{equation}
where \(a_k\in\{0,1\}\) is a parity determined by the \(T\)-matrix and \(\varepsilon_k=\pm1\) (see \cite{qefe}, Theorem~3.2).

\subsection{The Categorical Dirichlet Series and the QEFE}

The \textbf{categorical Dirichlet series} is defined for \(\Re(s)>1\) by
\begin{equation}\label{eq:Psi}
\Psi_k(s) = \sum_{j} \chi_{j,k}\, m_{j,k}^{-s}.
\end{equation}
It is a finite Dirichlet polynomial.
The Mellin transform of \(\Theta_k\) yields, for \(\Re(s)>1\),
\[
\pi^{-s/2}\,\Gamma\!\left(\frac{s+a_k}{2}\right) \Psi_k(s/2)\, D_k^{s}
= \int_0^\infty t^{s/2-1}\,\Theta_k(it)\,dt.
\]
Using the modular relation \eqref{eq:modrel} and splitting the integral gives the \textbf{Quantum‑Egyptian Functional Equation} (\cite{qefe}, Theorem~3.5):
\begin{equation}\label{eq:QEFE}
\Lambda_k(s) := \left(\frac{D_k^2}{\pi}\right)^{\!s/2} \Gamma\!\left(\frac{s+a_k}{2}\right) \Psi_k(s/2)
\end{equation}
satisfies
\begin{equation}\label{eq:functional}
\Lambda_k(s) = \varepsilon_k\,\overline{\Lambda_k(1-s)}.
\end{equation}
Equivalently,
\[
\Psi_k(s) = \varepsilon_k \left(\frac{D_k^2}{\pi}\right)^{\!s-\frac12}
\frac{\Gamma\!\left(\frac{1-s+a_k}{2}\right)}{\Gamma\!\left(\frac{s+a_k}{2}\right)}
\overline{\Psi_k(1-s)}.
\]
The completed function \(\Lambda_k(s)\) is an entire function of order~1, real on the real line, and its zeros are symmetric about \(\Re(s)=1/2\).
Because \(\Psi_k(s/2)\) is a finite polynomial, \(\Lambda_k(s)\) has finitely many zeros.

\subsection{The Limit to \(\zeta(2s)\)}

For the tower \(\mathcal{C}_k=\mathrm{SU}(2)_k\) one considers the normalised series
\begin{equation}\label{eq:Psihat}
\widehat{\Psi}_k(s) = \frac{\Psi_k(s)}{\Psi_k(1/2)}.
\end{equation}
Numerical computation (\cite{qefe}, Section~4) shows that as \(k\) increases, \(\widehat{\Psi}_k(s)\) approaches \(\zeta(2s)\) uniformly on compact subsets of \(\Re(s)>1/2\) (excluding the pole at \(s=1/2\)).
A heuristic explanation is that \(m_{j,k}\sim \text{const}\cdot (k+2)^3/(2j+1)^2\) and the Galois signs average to \(+1\), so the sum approximates \(\sum_{n\ge1} n^{-2s} = \zeta(2s)\).
The rigorous proof of this convergence is one of the central tasks of the programme.

\subsection{Explicit Formula for Finite Dirichlet Polynomials}

For a finite Dirichlet polynomial \(F(s)=\sum_{n\le N} a_n n^{-s}\) that satisfies a functional equation of the above type, there is an explicit formula linking its zeros to the coefficients.
The standard derivation integrates \(\Lambda_k'(s)/\Lambda_k(s)\) against a smooth even test function \(h\) with compactly supported Fourier transform.
The result is
\begin{equation}\label{eq:explicit}
\sum_{\rho_k} \widehat{h}\!\left(\frac{\rho_k-1/2}{2\pi i}\right)
= \widehat{h}\!\left(\frac{a_k}{2\pi i}\right) + \widehat{h}\!\left(\frac{-a_k}{2\pi i}\right)
- \sum_{j} \chi_{j,k}\, \frac{\log m_{j,k}}{m_{j,k}^{1/2}}\, h(\log m_{j,k}) + \text{gamma terms},
\end{equation}
where the sum over zeros \(\rho_k\) is finite.
This is the finite‑level analogue of the Weil explicit formula.

\section{Research Programme}

The programme is divided into five phases.
The ultimate goal is to prove that, in the limit \(k\to\infty\), all zeros of \(\Lambda_k(s)\) lie on the critical line \(\Re(s)=1/2\), thereby establishing the Riemann Hypothesis for \(\zeta(2s)\).

\subsection{Phase 1 --- Explicit Construction and Numerical Exploration}

\textbf{Objectives:} Compute \(\Psi_k(s)\) and its zeros for small to moderate \(k\); verify the functional equation; map the zero distributions; identify patterns and convergence towards the zeros of \(\zeta(2s)\).

\begin{enumerate}
\item Write code to compute \(m_{j,k}\) and \(\chi_{j,k}\) for all \(j\) and for \(k\le 500\) (or as high as computational resources allow).
  The Galois signs are determined from the action of a primitive root modulo \(4(k+2)\) on the \(S\)-matrix entries.
\item Form the Dirichlet polynomial \(\Psi_k(s)\) and the completed function \(\Lambda_k(s)\).
  Verify the functional equation numerically for random \(s\) and for real \(s\).
\item Compute the zeros of \(\Lambda_k(s)\) using standard root‑finding for entire functions (sign changes on the critical line, argument principle for off‑line zeros).
  Verify symmetry and count.
\item For increasing \(k\), plot the normalised zero distribution (histograms of ordinates scaled by \(\log k\)).
  Compare with the zeros of \(\zeta(2s)\) (from existing tables).
  Track the movement of individual zeros as \(k\) grows.
\item Formulate conjectures about the rate of convergence and the proportion of zeros already on the critical line for finite \(k\).
\end{enumerate}

\textbf{Deliverable:} A database of zeros for \(k\le 500\), plots showing convergence, and a precise conjecture about \(\sup_{\rho}|\Re(\rho)-1/2|\) as a function of \(k\).

\subsection{Phase 2 --- Rigorous Asymptotic Analysis}

\textbf{Objectives:} Prove \(\widehat{\Psi}_k(s)\to\zeta(2s)\) as \(k\to\infty\) uniformly on compact subsets of \(\Re(s)>1/2\) (away from the pole), with error bounds.

\begin{enumerate}
\item Write \(m_{j,k} = \frac{k+2}{4}\csc^2\bigl((2j+1)\frac{\pi}{k+2}\bigr)\) and use the Taylor expansion of \(\csc^2\) to obtain
  \[
  m_{j,k} = \frac{(k+2)^3}{3\pi^2(2j+1)^2} + O(k),
  \]
  with a uniform bound in \(j\).
\item Determine the Galois signs \(\chi_{j,k}\) explicitly (e.g.\ as a Jacobi symbol) and prove that for fixed \(j\) they are ultimately \(+1\) as \(k\to\infty\) along suitable subsequences.
  Prove that the average of \(\chi_{j,k}\) over a range of \(j\) tends to \(1\) sufficiently fast.
\item Approximate the sum \(\sum_j \chi_{j,k} m_{j,k}^{-s}\) by \(\sum_j (c_k(2j+1)^{-2})^{-s}\) with \(c_k = (k+2)^3/(3\pi^2)\).
  Show that the scaled sum \(\widehat{\Psi}_k(s)\) converges to \(\zeta(2s)\) uniformly on compact subsets of \(\Re(s)>1/2\).
\item Extend the convergence to \(\Re(s)>1/2-\delta\) using the functional equation, by comparing the completed functions \(\Lambda_k(s)\) and the completed Riemann zeta function \(\xi(2s)\).
\end{enumerate}

\textbf{Deliverable:} A theorem stating:
\textit{For any compact set \(K\subset\{s:\Re(s)>1/2\}\) not containing \(s=1/2\),
\[
\widehat{\Psi}_k(s) \longrightarrow \zeta(2s)
\]
uniformly, with an explicit error bound \(O(k^{-\alpha})\) for some \(\alpha>0\).}

\subsection{Phase 3 --- The Finite‑Level Explicit Formula}

\textbf{Objectives:} Derive an explicit formula for the zeros of \(\Lambda_k(s)\) and express the difference from the classical explicit formula for \(\zeta(2s)\) in terms of the theta function.

\begin{enumerate}
\item Starting from the Mellin representation, perform a contour integration to obtain the exact formula \eqref{eq:explicit}.
\item Write the analogous explicit formula for \(\zeta(2s)\) and subtract the two.
  The difference of the zero sums equals the difference of the weighted coefficient sums plus local terms.
\item Express the difference of the coefficient sums as an integral involving the twisted theta function \(\Theta_k\) and the classical theta function \(\Theta_\zeta\) (associated to \(\zeta(2s)\)).
\item Use the modular relation \eqref{eq:modrel} to decompose \(\Theta_k\) into a main term (matching the classical theta) and a remainder.
  Bound the remainder using Poisson summation and the lattice structure underlying \(\mathrm{SU}(2)_k\).
\end{enumerate}

\textbf{Deliverable:} An identity
\[
\sum_{\rho_k} \widehat{h}(\cdots) - \sum_{\rho_\zeta} \widehat{h}(\cdots) = \mathcal{E}_k(h),
\]
where \(\mathcal{E}_k(h)\) is expressed in terms of the modular tail of \(\Theta_k\) and is bounded by \(O(k^{-c})\) for suitable test functions.

\subsection{Phase 4 --- Zero Distribution and the Critical Line}

\textbf{Objectives:} Prove that as \(k\to\infty\), the proportion of zeros of \(\Lambda_k(s)\) not on the critical line tends to zero, or that the maximum real‑part deviation tends to zero.

\begin{enumerate}
\item Choose test functions whose Fourier transforms are supported in small intervals.
  The explicit formula from Phase~3 then shows that the zero‑counting measures \(\mu_k\) of \(\Lambda_k\) and \(\mu_\zeta\) of \(\zeta(2s)\) are close in the weak‑\(*\) topology.
\item Adapt Levinson’s mollifier method \cite{levinson} to the finite‑level setting.
  Multiply \(\Lambda_k(s)\) by a suitable Dirichlet polynomial \(M_k(s)\) to dampen the coefficients, and study the zeros of the mollified function.
  Use the exact functional equation to prove that a positive proportion of zeros lie on the critical line for each finite \(k\), and that this proportion tends to \(1\) as \(k\to\infty\).
\item Alternatively, use Rouché’s theorem to compare \(\Lambda_k(s)\) and \(\xi(2s)\).
  If the error between the two is uniformly small in a neighbourhood of a zero of \(\xi(2s)\), then the zeros are in one‑to‑one correspondence and close to each other.
  Assuming RH for \(\zeta(s)\) (as a target), the zeros of \(\Lambda_k\) must also lie on the line for large \(k\).
\item Prove the contrapositive: if a zero of \(\zeta(s)\) were off the critical line, it would force a zero of \(\Lambda_k(s)\) bounded away from the line for all large \(k\), contradicting the convergence \(\Lambda_k\to\xi(2s)\).
\end{enumerate}

\textbf{Deliverable:} A theorem:
\textit{If \(\widehat{\Psi}_k(s)\to\zeta(2s)\) uniformly on compact subsets of the critical strip with error \(O(k^{-c})\), then all zeros of \(\zeta(s)\) lie on the critical line.}

\subsection{Phase 5 --- Proof of the Riemann Hypothesis}

\textbf{Objectives:} Assemble the previous phases into a complete proof, or identify the exact obstacles that remain.

\begin{enumerate}
\item The main gap is likely the uniformity of convergence in the whole critical strip.
  Use Phragmén–Lindelöf and the modular transformation of the theta function to extend the convergence to \(\Re(s)=1/2\).
\item The error bound must be sharp enough to apply Rouché’s theorem.
  This will require analysing the theta series of the lattice underlying the MTC, which is related to classical Jacobi forms \cite{huang2019}.
\item If the full RH proof cannot be completed, extract partial results:
  \begin{itemize}
  \item A new zero‑free region for \(\zeta(s)\).
  \item A proof that 100\% of zeros lie on the critical line (a density theorem).
  \item A new proof of the Prime Number Theorem with an improved error term.
  \end{itemize}
\item Compare the categorical explicit formula with the Hilbert–Pólya idea.
  The zeros of \(\Lambda_k\) are eigenvalues of an operator related to the modular \(S\)-matrix; in the limit this could realise the spectral interpretation of the zeta zeros.
\end{enumerate}

\textbf{Deliverable:} Either a proof of the Riemann Hypothesis for \(\zeta(s)\), or a rigorous identification of the missing estimate together with a set of unconditional new results about the zeta zeros.

\section{Expected Milestones}

\begin{center}
\begin{tabular}{c p{10cm}}
\toprule
\textbf{Year} & \textbf{Milestone} \\
\midrule
1 & Numerical database of zeros for \(k\le 500\); conjectured convergence rate. \\
2 & Rigorous proof that \(\widehat{\Psi}_k(s)\to\zeta(2s)\) uniformly on compact subsets of \(\Re(s)>1/2\) with an explicit error term. \\
3 & Finite‑level explicit formula derived; error expressed in terms of the theta function; initial bounds on zero displacement. \\
4 & Proof that the proportion of zeros of \(\Lambda_k(s)\) on the critical line tends to \(1\) as \(k\to\infty\) (or that the maximal off‑line deviation tends to zero). \\
5 & Attempted proof of RH; or a new zero‑free region and a density theorem. \\
\bottomrule
\end{tabular}
\end{center}

\section{Concluding Remarks}

This programme avoids the pitfalls of the earlier sieve‑based attempts by working directly with the multiplicative and modular structure that the QEFE preserves.
The finite‑level Dirichlet polynomials \(\Psi_k(s)\) are natural objects from the representation theory of quantum groups and possess a genuine symmetry that mirrors the functional equation of \(\zeta(s)\).
Their limit to \(\zeta(2s)\) is a concrete, numerically testable phenomenon that, if rigorously established, would place the Riemann zeta function at the boundary of a family of exactly solvable spectral problems.
Even partial success would significantly advance our understanding of the zeta zeros.

\begin{thebibliography}{99}

\bibitem{greedy-continuation}
V.\ Geere, \emph{Greedy Harmonic Continuation}, 2026.\\
\url{https://victorgeere.co.za/maths/rh/research/greedy-harmonic/greedy-harmonic-continuation.html}

\bibitem{greedy-sieve}
V.\ Geere, \emph{The Greedy Harmonic Sieve and Classical Sieve Theory}, 2026.\\
\url{https://victorgeere.co.za/maths/rh/research/greedy-harmonic/greedy-harmonic-sieve.html}

\bibitem{qefe}
V.\ Geere, \emph{A Quantum Egyptian Functional Equation}, v8, 2026.\\
\url{https://victorgeere.co.za/maths/rh/research/qefe/quantum-egyptian-functional-equation-v8.html}

\bibitem{mds}
Author, \emph{The Multiplicative Depth Sieve}, preprint, 2025.

\bibitem{weil}
A.\ Weil, \emph{Sur les formules explicites de la th\'eorie des nombres}, Comm.\ S\'em.\ Math.\ Univ.\ Lund, tome suppl.\ d\'edi\'e \`a M.\ Riesz (1952), 252--265.

\bibitem{huang2019}
Y.-Z.\ Huang, \emph{Riemann hypothesis, modular tensor categories and vertex operator algebras}, 2019.

\bibitem{levinson}
N.\ Levinson, \emph{More than one third of zeros of Riemann's zeta-function are on \(\sigma=1/2\)}, Adv.\ Math.\ 13 (1974), 383--436.

\bibitem{iwaniec}
H.\ Iwaniec \& E.\ Kowalski, \emph{Analytic Number Theory}, AMS Colloq.\ Publ.\ 53, 2004.

\bibitem{titchmarsh}
E.C.\ Titchmarsh, \emph{The Theory of the Riemann Zeta-function}, 2nd ed., Oxford, 1986.

\end{thebibliography}

\end{document}